% p_entries.tex - Terms beginning with P
\chapter*{P}
\addcontentsline{toc}{chapter}{P}

\dictentry{PIN}{/pɪn/, \textit{abbr.}}{
Personal Identification Number. A numeric password used to authenticate a user
to a system, most commonly for debit or credit card transactions at an ATM or
POS terminal. PINs are highly sensitive and are typically protected by
encryption (like DUKPT) immediately upon entry.

\example{When you buy groceries and the terminal asks for your "code," you are
entering your PIN to authorize the transaction.}

\analogy{Biological}{Like a \textbf{Receptor Binding Motif}. Just as a specific
molecule (like a hormone) must have a precise shape to fit into a cellular
receptor and "authorize" a biological response, your PIN must have the precise
sequence to "fit" the bank's security check and authorize your payment.}

\crossref{\textbf{Authentication}, \textbf{DUKPT}, \textbf{HSM}}

\etymology{Introduced with the first ATMs in the late 1960s. The concept was
patented by James Goodfellow in 1966.}
}
\indexentry{PIN}

\dictentry{Proximity Card}{/prɒkˈsɪmɪti kɑːrd/, \textit{n.}}{
A contactless smart card that must be held very close to a reader—usually
within 10 centimeters (4 inches)—to communicate. This is the technology used
in most modern contactless credit cards and mobile wallets.

\crossref{\textbf{ISO/IEC 14443}, \textbf{Vicinity Card}, \textbf{NFC}}
}
\indexentry{Proximity Card}

% Additional P terms will be added here
